\subsection{A boundary root approaching the central line}
\label{sec:weak-root-passage}

When the constant coefficient is small compared with the linear one,
the first boundary root on the positive branch approaches the central
line relative to the current coefficient scale. Rescaling by its radius
keeps this root simple and near \(x=1\). In the resulting logarithmic
height coordinate, the positive endpoint has a large negative
logarithmic rate. A uniform lower bound for the negative endpoint's
rate then makes \(\mathcal L\) increase throughout the root interval.

Let \(a,\kappa>0\) be small and \(\hat r\in(0,1]\), and suppose
\begin{equation}
\label{eq:weak-root-scales}
 q_0=-\kappa\ell^2,\qquad q_1=\kappa\ell,\qquad q_2=-\epsilon_2\kappa,
 \qquad \ell=\hat r\kappa a^2,\qquad |\epsilon_2|\le C_2a,
\end{equation}
with a fixed \(C_2\). In \(|\eta|=\ell x\), \(z=\ell^2Y\) the outward
drifts and height equations are
\[
 K_\sigma(x)=-\frac{q(\sigma\ell x)}{\kappa\ell^2}
 =1-\sigma x+\epsilon_2x^2-\sigma\hat ra^2x^3R_q(\sigma\ell x),\qquad
 x_{\log Y}=\kappa K_\sigma(x)/x+YB(\sigma\ell x,\ell^2Y)/x .
\]
For a fixed small preliminary upper bound on \(a\), the first positive
root of \(K_+\) satisfies \(\alpha=1+O(a)\) and \(K_+'(\alpha)<-c<0\).
Put \(x_L=\alpha/4\), \(w_a=\alpha/2\) and \(R_p=3\alpha/4\), and use
the logarithmic height coordinates \(T=-\kappa a\log h\) and
\(\widehat T=T/a=-\kappa\log h\).
We compare the branches in \(T\), after estimating the positive branch
in \(\widehat T\).
For a fixed small height \(Y_f>0\), the root interval is
\begin{equation}
\label{eq:weak-root-interval}
 w=\mathcal W_+(x(\ell^2Y_f;h),Y_f)\in[w_a,\alpha).
\end{equation}

The negative half-orbit need not be at the weak-root radius when the
positive half-orbit approaches its root. To estimate it on the same
central heights, introduce a variable radius scale: let \(d>0\)
and put \(\ell_d=\ell d/a\), \(\delta_d=\kappa a/d\) and
\(a_d=a/d\), so that \(a_d/\delta_d=1/\kappa\).
In \(|\eta|=\ell_dx\), the outward drift on the negative branch is
\[
 K_d(x)=a_d+x+(\epsilon_2/a)dx^2+\hat rd^2x^3R_q(-\ell_dx),
\]
so that \(K_d(x)=K_1(dx)/d\).
Fix \(d_{\max}>0\) such that, for \(0<d\le d_{\max}\),
\begin{equation}
\label{eq:weak-root-negative-drift}
 K_d(x)\asymp a_d+x,\qquad K_d'(x)\ge\tfrac12\qquad(0\le x\le4),
\end{equation}
with constants depending on \(d_{\max}\). This holds for every small
\(d_{\max}\), and for every fixed \(d_{\max}\) when \(\epsilon_2\ge0\) and
\(\hat r\) is small.

\begin{lemma}[Two branches at the weak root]
\label{lem:weak-root-branches}
There are \(c_1,c_d>0\) and a fixed small \(Y_b>0\) such that, for
sufficiently small \(a\) and \(\kappa\), the following hold in the coordinate
\(T\).
\begin{enumerate}[label=\textup{(\roman*)}]
\item On the root interval \eqref{eq:weak-root-interval},
\(\mathfrak H_+\le-c_1/a\).
\item Let \(c_da\le d\le d_{\max}\), and let \(J\subset[1/2,x_{\max}]\) be
a compact interval, with a fixed \(x_{\max}\), such that \(K_d>0\) on
\((0,\max J)\). At the height \(\ell_d^2Y_b\), and for negative projected
states in \(J\),
\[
 d\,\mathfrak H_-\ge\min_{x\in J}\frac{K_d'(x)}x-o(1),\qquad
 |I_{-,T}|\le CE_d/d=o(1),\qquad
 E_d=\ell_d^2\bigl(z_*/(\ell_d^2Y_b)\bigr)^{-2q_2},
\]
with \(o(1)\to0\) as \(\kappa\to0\) uniformly in \(d\) and \(J\), up to a
root of \(K_d\) at \(\max J\); the constants and the smallness of
\(a,\kappa\) depend on \(x_{\max}\). In particular,
\(\mathfrak H_-\ge c/d>0\) for \(J=[0.9,4]\).
\item On \eqref{eq:weak-root-interval}, the negative graph has
\(z<(Y_b/2)|\eta|^2\) at \(|\eta|=\ell c_d\).
\end{enumerate}
\end{lemma}

\begin{proof}
(i) The boundary-flow coordinate applies with small parameter
\(\kappa\), since the error from \(B\) is \(O(\ell)=O(\kappa)\). At
\(t=\ell^2Y_f\) the coefficient ratios are \(|q_0|/t=\kappa/Y_f\),
\(q_1/\sqrt t=\kappa/\sqrt{Y_f}\) and \(|q_2|\le C_2\kappa a\). Hence
\cref{lem:boundary-flow-endpoint} gives
\(-I_{+,w}=E_aA_+wV_+(w)\), with
\(E_a=\ell^2(z_*/(\ell^2Y_f))^{-2q_2}\) and
\(V_+=1+O_{C^3}(\kappa)\). The fixed positive phase advance from
\(x_L\) to \(w_a\) allows us to apply
\cref{lem:direct-weighted-phase} in the coordinate \(\widehat T\).
Since \(K_+'(w)/w\) has a fixed negative margin on
\([w_a,\alpha)\), \cref{cor:direct-weak-root-comparison} gives (i).

(ii) By \eqref{eq:weak-root-negative-drift} and \(a_d/\delta_d=1/\kappa\ge1\),
the case of a small central coefficient in \cref{lem:direct-source-height},
with \(K_*=a_d\), gives a uniform height unit at the entry \(x_i=1/4\),
also as \(a_d\to0\).
The boundary-flow projection at \(Y_b\), on a fixed radius interval
extending \([1/4,x_{\max}]\), has inverse energy error
\(O_{C^k}(\delta_dY)\), since \(\ell_d/\delta_d=\hat rd^2\) is bounded.
By \cref{lem:direct-boundary-flow}, the projection's height bound \(Y_0\) depends
only on the lower end of that interval, so \(Y_b\) does not depend on
\(x_{\max}\).

For output states in \(J\), the advance from \(x_i\) to \(1/2\)
takes a fixed positive boundary time, since \(K_d/x\) is bounded
there. Thus \cref{lem:direct-weighted-phase} applies uniformly up to
a root of \(K_d\). At \(\ell_d^2Y_b\), the coefficient ratios are
\(\delta_da_d/Y_b\), \(\delta_d/\sqrt{Y_b}\) and
\(|q_2|\le C_2\kappa a\), and \(\ell_d\to0\).
These quantities are small because \(\delta_d\le\kappa/c_d\)
and \(a_d\le1/c_d\). \Cref{lem:boundary-flow-endpoint} with
\(\ell=\ell_d\), \(\delta=\delta_d\) and \(J_w=[1/2,x_{\max}]\)
therefore gives \(-I_{-,w}=E_dA_dw(1+O_{C^3}(\delta_d))\).

To express these estimates in the common coordinate, use
\(\widehat T_d=-\delta_d\log h=T/d\), so
\(\partial_T=d^{-1}\partial_{\widehat T_d}\). In \(\widehat T_d\),
the first two lines of \eqref{eq:direct-rate-identities}, with
\(\theta-1,\theta'=o(1)\), give \(|I_{-,\widehat T_d}|\le CE_d\) and
\(\partial_{\widehat T_d}\log(hI_{-,h})\ge\min_JK_d'(x)/x-o(1)\).
Division by \(d\) gives the two estimates. On \([0.9,4]\),
\(K_d'(x)/x\ge1/8\) by \eqref{eq:weak-root-negative-drift}.

Finally, \(|q_2|\le C_2\kappa a<1/8\), so
\eqref{eq:direct-tail-amplitude} with \(t=\ell_d^2Y_b\) gives
\(E_d\le C\ell_d^{3/2}\) and
\(E_d/d\le C(\ell/a)^{3/2}d^{1/2}\to0\), uniformly for
\(0<d\le d_{\max}\), since \(\ell/a=\hat r\kappa a\to0\).

(iii) Choose \(c_d\) so small that the negative prefix action to
\(x=c_d\) in the weak-root scale is less than one quarter of the fixed
positive phase advance \(\tau_*\) from \(x_L\) to \(w_a\). On
\eqref{eq:weak-root-interval} this advance gives
\(h/(\kappa\ell^2)\le(Y_f/\kappa)\e^{-\tau_*/\kappa}\). The negative
central-height coordinate then gives
\[
 \frac{z(-\ell c_d;h)}{(\ell c_d)^2}
 \le\frac{CY_f}{c_d^2}\,\e^{-3\tau_*/(4\kappa)}
 <\frac{Y_b}{2}
\]
for small \(\kappa\): the exponential decay in \(1/\kappa\) absorbs the
bounded height units and the polynomial prefactors.
\end{proof}

\begin{corollary}[Negative branch up to its first root]
\label{cor:weak-root-negative-first-root}
Assume \eqref{eq:weak-root-scales} and \eqref{eq:weak-root-negative-drift},
with \(Y_b\) as in \cref{lem:weak-root-branches}, and suppose that the
negative drift \(K_1\) in \(|\eta|=\ell X/a\) has a first positive root
\(\beta\le\beta^*\), with a fixed \(\beta^*\). For sufficiently small \(a\)
and \(\kappa\), every central height whose negative graph has
\(z<Y_b|\eta|^2\) at \(|\eta|=\ell_{d_{\max}}\) satisfies, in the
coordinate \(T\),
\[
 \mathfrak H_-\ge\min_{d_{\max}/2\le X\le\beta}\frac{K_1'(X)}X-o(1),
 \qquad |I_{-,T}|=o(1),
\]
where \(o(1)\to0\) as \(\kappa\to0\), uniformly as the negative projected
state at the height \(\ell_{d_{\max}}^2Y_b\) approaches the root.
\end{corollary}

\begin{proof}
By \eqref{eq:weak-root-negative-drift} at \(d=d_{\max}\) and
\(K_d(x)=K_1(dx)/d\), \(K_1>0\) on \((0,4d_{\max}]\).
The first positive root of \(K_{d_{\max}}\) is therefore
\(\beta/d_{\max}>4\). In \(|\eta|=\ell_{d_{\max}}x\),
\(z=\ell_{d_{\max}}^2Y\), the graph reaches \(x=1\) below
\(Y=Y_b\). Equation~\eqref{eq:direct-projection-energy} puts its
projected state there above \(1/2\), since \(Y_b\) and \(\kappa\)
are small. The projected state then increases; by
\cref{rem:direct-projected-continuation}, it does not cross
\(\beta/d_{\max}\), and the graph continues to \(Y_b\).

At that height the projected state lies in
\(J=[1/2,\beta/d_{\max}]\). Apply part~(ii) of
\cref{lem:weak-root-branches} with \(d=d_{\max}\), this \(J\), and
\(x_{\max}=\beta^*/d_{\max}\). The assertion follows from
\(K_d'(x)/x=d\,K_1'(dx)/(dx)\), with \(d_{\max}\) fixed.
\end{proof}

On the root interval, a negative graph is controlled by
\cref{lem:weak-root-branches} if \(z/|\eta|^2\) reaches \(Y_b\) in the
prescribed radius range. The following proposition isolates the
additional estimate needed for graphs that remain below this level.
Each application will verify that estimate on its own coefficient
scale.

\begin{proposition}[Weak root]
\label{prop:weak-root-passage}
Assume \eqref{eq:weak-root-scales} and \eqref{eq:weak-root-negative-drift},
and let \(Y_b\) be as in \cref{lem:weak-root-branches}. Suppose there are
\(C_n\) and \(\varepsilon_p\to0\) such that every central height in the
root interval \eqref{eq:weak-root-interval} whose negative graph has
\(z<Y_b|\eta|^2\) at \(|\eta|=\ell_{d_{\max}}\) satisfies
\begin{equation}
\label{eq:weak-root-outer}
 \mathfrak H_-\ge-C_n,\qquad |I_{-,T}|\le\varepsilon_p .
\end{equation}
Then there are \(a_\circ,\kappa_\circ>0\) such that, for \(0<a\le a_\circ\),
\(0<\kappa\le\kappa_\circ\) and all sufficiently small remaining parameters,
every zero of \(\mathcal L\) on \(\cH_p\) satisfies
\(\mathcal L_T>0\). Consequently \(\Delta_p\) has at most two distinct
zeros there.
\end{proposition}

\begin{proof}
\emph{The negative branch on the root interval.} By
\cref{lem:weak-root-branches}(iii), the ratio \(z/|\eta|^2\) along the
negative graph is below \(Y_b/2\) at \(|\eta|=\ell c_d\).
If the ratio reaches \(Y_b\) at a
radius \(\ell_d\) with \(c_da<d\le d_{\max}\), choose the first
such \(d\) for the central height under consideration. At
\(\ell_d^2Y_b\), the normalized radius is then one, and small
\(Y_b\) puts the projected state in \([0.9,4]\). Hold this section
fixed when varying the central height locally. The resulting
representation of \(I_-\) agrees with \(I_-(h)\), so
\cref{lem:weak-root-branches}(ii) applies as in
\cref{lem:direct-frozen-section-sign}. Otherwise, the ratio is still
below \(Y_b\) at \(|\eta|=\ell_{d_{\max}}\), and
\eqref{eq:weak-root-outer} applies. In both cases,
\(\mathfrak H_-\ge-C_n\) and \(|I_{-,T}|=o(1)\) throughout
\eqref{eq:weak-root-interval}.

\emph{The root interval.} Combine this bound for the negative branch
with \cref{lem:weak-root-branches}(i) for the positive branch.
\Cref{cor:direct-weak-root-comparison} gives \(\mathcal L_T>0\)
throughout \eqref{eq:weak-root-interval}, after \(a_\circ\) is reduced.
The constants are fixed on the preliminary small-\(a\) region.

\emph{The height cover.} Apply \cref{lem:prefix-root-cover} with
\(\delta=\kappa\), \(\epsilon_3=\ell/\kappa=\hat ra^2\), comparison drifts
\(1-\sigma x+\epsilon_2x^2\), and the radii \(w_a,R_p,\alpha\) above.
The comparison drifts have the required uniform positive margin, the
gap \(R_p-w_a\) is uniform, and the boundary-flow inverse is defined
up to \(\alpha\) at \(Y_f\). Together with \(\mathcal L_T>0\) on
the root interval, the lemma gives at most two distinct zeros of
\(\Delta_p\) on \(\cH_p\).
\end{proof}
